\documentclass[10pt,a4paper]{amsart}
\usepackage[margin=19mm]{geometry}
\usepackage{amsmath,amssymb,mathtools}
\usepackage[scr=rsfs,cal=euler]{mathalfa}
\usepackage{xfrac}
\usepackage[unicode,hidelinks,bookmarks=false]{hyperref}
\newcommand{\ie}{i.e.\ }
\newcommand{\cf}{cf.\ }
\newcommand{\resp}{respectively\ }
\newcommand{\Subsection}{\subsection}
\providecommand{\nobreakdash}{\nobreak\hbox{-}}
\setlength{\emergencystretch}{2em}
\multlinegap0pt
\usepackage{amssymb}
\usepackage{xfrac}
\usepackage[shortlabels]{enumitem}
\usepackage{dsfont}
\usepackage{mathtools}
\newtheorem{prop}{Proposition}
\newtheorem{lemma}{Lemma}
\newcommand{\SR}{{\sf SR}}
\newcommand{\eset}{\varnothing}
\title{Cutoff with an $O(1)$ window for Potts Glauber Dynamics\\on lattice at High Temperature}
\author{Seoyeon Yang, Allan Sly}
\date{Source: arXiv:2608.26259v1 (CC BY 4.0)}
\begin{document}
\maketitle

\noindent\emph{Source and licence.} Cutoff with an $O(1)$ window for Potts Glauber Dynamics\\on lattice at High Temperature. Version arXiv:2608.26259v1 (CC BY 4.0), distributed by arXiv under the Creative Commons Attribution 4.0 International licence, http://creativecommons.org/licenses/by/4.0/, as recorded in the arXiv licence metadata for this exact version and shown on the abstract page https://arxiv.org/abs/2608.26259v1. Full licence notice: \texttt{licenses/arxiv-2608.26259-CC-BY-4.0.txt}.\par\smallskip

\noindent\textit{Editorial scope.} Counted unit: the article's prop prop:one-block-SR (source0.tex lines 5384--5402). The article's complete original proof of it is delivered with it (source0.tex lines 6906--7268, one \texttt{proof} environment of 363 lines, which the article places away from the statement). No mathematics is written, repaired or re-derived: the statement and the proof are the article's own lines, in the article's own notation, and no retained line is substituted or re-flowed.  Retained uncounted as the source-local context the proof needs: lines 4441--4443; lines 4642--4651; lines 6113--6121; lines 6502--6523; lines 6633--6663.  Nothing was omitted from any retained span, so the package carries no omission record.  No retained line contains a citation, so the wrapper carries no bibliography and no bibliography entry.  No retained line contains a figure environment, an image-inclusion command or drawing code, so no image is delivered and the package declares no asset dependency.  Editorial formatting only, in addition: an AMS \texttt{amsart} preamble that restates the article's own declarations and macros used by the retained lines, namely the environments \texttt{prop}, \texttt{lemma} on separate counters, together with the standard packages those lines need.  The retained environments print in this document as Proposition 1, Lemma 1 in order of appearance, whereas the article prints the same items with the numbering of its own full document.  The complete native source of the article as distributed in the arXiv e-print for this exact version is bundled unchanged as \texttt{sources/208624/source0.tex}.
\begin{equation}\label{eq:cone-embedded}
2\bigl(\lceil c_{\rm cone}T\rceil+1\bigr)<n.
\end{equation}

\begin{prop}[Residual weight of an isolated cone]\label{prop:isolated-cone-residual}
Under \eqref{eq:cone-embedded}, there exists \(C<\infty\) such that, for every
initial configuration \(\sigma_0\in[q]^{\Lambda_n}\) and every
\((v,t)\in\Lambda_n\times[0,T]\),
\[
\mathbb{P}_{\sigma_0}\left(
\mathsf{Bot}_t(v),\,J_{(v,t)}>p_t(v)
\right)\leq C e^{-(1+\kappa)t}.
\]
\end{prop}

\begin{equation}\label{eq:top-unit-U}
\mathcal U_A=\mathcal U_A(X)
        :=
        \bigcap_{u\in A_{\mathrm{top}}}
        \left\{
        \text{there is at least one update at }u
        \text{ in }(s_u,T]
        \right\}.
\end{equation}

\begin{lemma}[First admissible cone estimate]
\label{lem:first-h-admissible}
Fix \(\lambda_0>0\). After decreasing
\(\beta_0(d,q)\), there are constants \(C,c<\infty\) such that the following
holds.

Fix \(h:=\ell_A(X)\). On \(\mathcal R_1\), let \(r_{A,1}\in[h,T]\) be the
first \((A,h)\)-admissible backward depth in the sequential exploration above. Let
\(\mathcal N_0\) be the event that \(A\) is Red and \(\mathcal R_1^c\) occurs.
\[
        \mathbb E
        \left[
            e^{(1+\kappa)r_{A,1}}\mathbf 1_{\mathcal R_1}
            +
            e^{(1+\kappa)T}\mathbf 1_{\mathcal N_0}
            \,\middle|\,
            \mathcal U_A
        \right]
        \le
        C e^{-\lambda_0 W(A)}e^{c\beta h}.
\]
\end{lemma}

\begin{lemma}[Weighted continuation bound]
\label{lem:one-site-continuation}
After decreasing \(\beta_0(d,q)\), there exists \(C_{\rm cont}<\infty\), depending only on \(d\) and \(q\), such that the following holds.

Fix a finite nonempty \(A\subset\Lambda\), \(h\in[0,T]\), and a reached
\((A,h)\)-admissible trial at depth \(r\).  Let \(v=v(r)\), and let
\(\mathscr G\) be the sigma-field generated by the exploration revealed before
testing the cone at \((v,T-r)\), including the data
\((r,v,\mathsf C_r(v))\), but excluding the graphical marks strictly below its tip. 
Continue the sequential exploration from this trial, and write
\[
        r=r_0<r_1<\cdots<r_J
\]
for the depths of all trials subsequently reached, including the current one.
Let
\[
        \mathcal N_r
        :=
        \{\text{the continuation reaches depth \(T\) without a successful trial}\}.
\]
Then
\begin{equation}\label{eq:one-site-continuation}
\mathbb E\left[
    \sum_{j=0}^{J}e^{(1+\kappa)(r_j-r)}
    +e^{(1+\kappa)(T-r)}\mathbf 1_{\mathcal N_r}
    \,\middle|\,
    \mathscr G
\right]
\le C_{\rm cont}.
\end{equation}
\end{lemma}

\begin{prop}
\label{prop:one-block-SR}
For every \(\lambda_{\rm SR}>0\), after decreasing \(\beta_0(d,q)\) if necessary, there exist constants \(C_{\mathrm{sr}}<\infty\) and \(c_{\mathrm{sr}}<\infty\) such that the following holds.  For every non-empty \(A\subseteq\Lambda\), and \(0\le h\le T \),
\[
    \sup_{\substack{X:\ \ell_A(X)=h}}
    q^{|A|}
    \mathbb P\left(
        A\in\mathrm{SR}
        \,\middle|\,
        \mathcal H_A^{-}=X,
        \mathsf T_A
    \right)
    \le
    C_{\mathrm{sr}}
    e^{-(1+\kappa)T}
    e^{-\lambda_{\rm SR} W(A)}
    e^{c_{\mathrm{sr}}\beta h}.
\]
\end{prop}

\begin{proof}[Proof of Proposition~\ref{prop:one-block-SR}]


We suppress the copy index \(a\) throughout the proof.
Fix the outside transcript \(\mathcal{H}_A^-=X\), and set
\[
        h:=\ell_A(X).
\]


%%%%
For a given subset $S\subset \Lambda$, define $V_{\mathrm{Blue}_S^*}$ be the union of the top sets of the Blue clusters that arise when
exposing the joint histories of $S$. Then,
\[
\{A\subset V_{\rm{Blue}}\}=\{A\subset V_{\mathrm{Blue}_A^*}\}\cap \{\mathcal{H}_A\cap X=\eset\}.
\]

%%%%
In an exploration from \(A\times\{T\}\), define
\[
        E_A^*(h):=E_{{\rm res},A}^*(h)\cup E_{{\rm fail},A}^*(h),
\]
where \(E_{{\rm res},A}^*(h)\) is the event that, while exposing only histories
of \(A\), some reached admissible cone trial \(k\) succeeds and takes the
residual branch
\[
        J_{(v_{A,k},\tau_{A,k})}>p_{\tau_{A,k}}(v_{A,k}),
\]
and \(E_{{\rm fail},A}^*(h)\) is the event that, while exposing only histories
of \(A\), the set \(A\) is Red but no admissible cone trial succeeds.
Then,
\[
\{A\in\SR\}\subset E_A^*(h)\cap\{\mathcal H_A\cap X=\eset\}.
\]
Remark that as $X\cap \mathsf D_h(A)=\varnothing$, every cone \(\mathsf C\subset\mathsf D_h(A)\) is disjoint from the outside
transcript \(X\). 


%%%
We use the top-unit event from \eqref{eq:top-unit-U} on time slab $[T-1,T]$.
On $\mathsf T_A:=\{A\in{\rm SR}\}\cup \{A\subset V_{\rm Blue}\}$, \(\mathcal U_A\) must occur.
Indeed, if some
\(u\in A\) has no update in \((s_u,T]\), then the vertical branch
from \((u,T)\) hits the outside transcript at \((u,s_u)\). Therefore,
$\mathbb{P}\big(A\in\mathrm{SR} \,\big\vert\, \mathcal{H}_A^-= X,\{A\in\mathrm{SR}\}\cup\{A\subset V_{\mathrm{Blue}}\}\big)$ is bounded above by



\begin{align*}
\frac{
\mathbb{P}\left(
E_A^*(h),\,
\mathcal{H}_{A}\cap  X=\emptyset,\,
\mathcal{U}_A
\mid \mathcal{H}_{A}^{-}= X
\right)
}{
\mathbb{P}\left(
A\subset V_{\mathrm{Blue}_{A}^{*}},\,
\mathcal{H}_{A}\cap X=\emptyset,\,
\mathcal{U}_A
\mid \mathcal{H}_{A}^{-}= X
\right)
}
=
\frac{
\mathbb{P}\left(
E_A^*(h),\,
\mathcal{H}_{A}\cap  X=\emptyset
\mid \mathcal{U}_A
\right)
}{
\mathbb{P}\left(
A\subset V_{\mathrm{Blue}_{A}^{*}},\,
\mathcal{H}_{A}\cap X=\emptyset
\mid \mathcal{U}_A
\right)
}
.
\end{align*}
The equality holds because all indicated events are now $\mathcal{H}_A$ measurable.




%%%
For the denominator, there is a constant
\(c_{\rm bl}>0\) such that
\[
\mathbb{P}\left(
A\subset V_{\mathrm{Blue}_{A}^{*}},\,
\mathcal{H}_{A}\cap X=\emptyset
\mid \mathcal{U}_A
\right)
        \ge
        c_{\rm bl}^{|A|}.
\]
Indeed, for \(u\in A_{\rm top}\), conditional on the existence of an update in
\((s_u,T]\), the latest such update is oblivious with probability
\(p_{\rm obl}\). For \(u\in A\setminus A_{\rm top}\), require that the latest
update in \((T-1,T]\) exists and is oblivious; this has probability
\((1-e^{-1})p_{\rm obl}\). If all these independent events occur, then the
histories from all sites of \(A\) die before time \(T-1\), and \(A\subset
V_{\rm Blue}\). Thus one may take $c_{\rm bl}:=p_{\rm obl}\min\{1,1-e^{-1}\}$. Therefore,
\begin{align*}
\mathbb P
\left(
    A\in{\rm SR}
    \,\middle|\,
    \mathcal H_A^-=X,
    \{A\subset V_{\rm Blue}\}\cup\{A\in{\rm SR}\}
\right)  
\leq
c_{\rm bl}^{-|A|}
        \mathbb P
        \left(
            E_A^*(h),\,
\mathcal{H}_{A}\cap  X=\emptyset
            \,\middle|\,
            \mathcal U_A
        \right)
\leq
c_{\rm bl}^{-|A|}
        \mathbb P
        \left(
            E_A^*(h)
            \,\middle|\,
            \mathcal U_A
        \right).
\end{align*}





It remains to bound $\mathbb P\left(E_A^*(h)\,\middle|\,\mathcal U_A\right)$.
Recall that \(\mathcal R_k\) is the event that trial \(k\) is reached, that
\(r_{A,k}\) is its backward depth on this event, and that
\(\tau_{A,k}=T-r_{A,k}\) is its absolute tip time, hence the remaining
lower-cone time.
Define
\[
        \mathcal M_A
        :=
        \sum_{k\ge1}
        \mathbf 1_{\mathcal R_k}
        e^{(1+\kappa)r_{A,k}}
        +
        e^{(1+\kappa) T}
        \mathbf 1_{E_{{\rm fail},A}^*(h)} .
\]
We first prove that, for every \(\lambda_0>0\), after decreasing
\(\beta_0(d,q)\) if necessary,
\begin{equation}\label{eq:MA-bound}
        \mathbb E\left[
            \mathcal M_A
            \,\middle|\,
            \mathcal U_A
        \right]
        \le
        C e^{-\lambda_0W(A)}e^{c\beta h}.
\end{equation}
On \(\mathcal R_1\), define the continuation functional from the first trial by
\[
\begin{aligned}
        \mathcal C_{A,1}
        &:=
        \sum_{j\ge0}
        \mathbf 1_{\mathcal R_{1+j}}
        e^{(1+\kappa)(r_{A,1+j}-r_{A,1})}
        +
        e^{(1+\kappa)(T-r_{A,1})}
        \mathbf 1_{E_{{\rm fail},A}^*(h)} .
\end{aligned}
\]
Then, pathwise on \(\mathcal R_1\), $\mathcal M_A
        =
        e^{(1+\kappa)r_{A,1}}\mathcal C_{A,1}$.
Let \(\mathcal G_{A,1}\) be the intrinsic pre-test sigma-field of the first
trial, enlarged by \(\sigma(\mathcal U_A)\). Since \(r_{A,1}\ge h\) on
\(\mathcal R_1\), and since
\(h=\ell_A(X)\ge1\) by the definition of \(\ell_A\), the top-unit event
\(\mathcal U_A\) is determined by the history exposed above the first trial.
Thus, on \(\mathcal R_1\), adding \(\mathcal U_A\) does not change the
pre-test state or the fresh lower-cone law.
Lemma~\ref{lem:one-site-continuation}, applied with
\(r=r_{A,1}\), gives
\[
        \mathbb E\left[
            \mathcal C_{A,1}
            \,\middle|\,
            \mathcal G_{A,1}
        \right]
        \le C_{\rm cont}.
\]
Since
\(e^{(1+\kappa)r_{A,1}}\mathbf 1_{\mathcal R_1}\) is
\(\mathcal G_{A,1}\)-measurable,
\[
\begin{aligned}
        \mathbb E\left[
            \mathcal M_A\mathbf 1_{\mathcal R_1}
            \,\middle|\,
            \mathcal U_A
        \right]
        &\le
        C_{\rm cont}
        \mathbb E\left[
            e^{(1+\kappa)r_{A,1}}
            \mathbf 1_{\mathcal R_1}
            \,\middle|\,
            \mathcal U_A
        \right].
\end{aligned}
\]
On \(\mathcal R_1^c\), no trial is reached. Hence
\[
        E_{{\rm fail},A}^*(h)\cap\mathcal R_1^c
        =
        \mathcal N_0,
\]
where \(\mathcal N_0\) is the no-admissible-trial Red event from
Lemma~\ref{lem:first-h-admissible}. Therefore
\[
\begin{aligned}
        \mathbb E\left[
            \mathcal M_A
            \,\middle|\,
            \mathcal U_A
        \right]
        &\le
        C
        \mathbb E\left[
            e^{(1+\kappa)r_{A,1}}
            \mathbf 1_{\mathcal R_1}
            +
            e^{(1+\kappa) T}\mathbf 1_{\mathcal N_0}
            \,\middle|\,
            \mathcal U_A
        \right] \le
        C e^{-\lambda_0W(A)}e^{c\beta h},
\end{aligned}
\]
where the last inequality is Lemma~\ref{lem:first-h-admissible}. This proves
\eqref{eq:MA-bound}.







We now convert the weighted estimate into a probability bound for
\(E_A^*(h)\). By definition,
\[
        E_{{\rm res},A}^*(h)
        =
        \bigcup_{k\ge1}
        \left\{
            \mathcal R_k,\,
            \mathsf{Bot}_{\tau_{A,k}}(v_{A,k}),\,
            J_{(v_{A,k},\tau_{A,k})}>p_{\tau_{A,k}}(v_{A,k})
        \right\}.
\]
Let \(\mathcal G_{A,k}\) be the \(k\)-th pre-test sigma-field enlarged
by \(\sigma(\mathcal U_A)\). On \(\mathcal R_k\), the lower cone below
\((v_{A,k},\tau_{A,k})\) is contained in \(\mathsf D_h(A)\), and
therefore is disjoint from the outside transcript \(X\). Since
\(r_{A,k}\ge h\ge1\), the event \(\mathcal U_A\) is already
determined by the pre-test history, and the lower cone is fresh conditionally
on \(\mathcal G_{A,k}\). Proposition \ref{prop:isolated-cone-residual}
applies when \(\tau_{A,k}>0\), while the case \(\tau_{A,k}=0\) is trivial after
increasing \(C\). Thus, in both cases, the following holds on \(\mathcal R_k\):
\[
\begin{aligned}
        \mathbb P
        \left(
            \mathsf{Bot}_{\tau_{A,k}}(v_{A,k}),
            J_{(v_{A,k},\tau_{A,k})}>p_{\tau_{A,k}}(v_{A,k})
            \,\middle|\,
            \mathcal G_{A,k}
        \right)
        &=
        M_{\tau_{A,k}}\bigl(1-p_{\tau_{A,k}}(v_{A,k})\bigr)
        \le
        C e^{-(1+\kappa)\tau_{A,k}}     
        =
        C e^{-(1+\kappa) T}e^{(1+\kappa)r_{A,k}} .
\end{aligned}
\]
The event \(\mathcal R_k\) is \(\mathcal G_{A,k}\)-measurable, so the
tower property and a union bound imply
\[
\begin{aligned}
        \mathbb P(E_{{\rm res},A}^*(h)\mid\mathcal U_A)
        &\le
        C e^{-(1+\kappa) T}
        \mathbb E
        \left[
            \sum_{k\ge1}
            \mathbf 1_{\mathcal R_k}
            e^{(1+\kappa)r_{A,k}}
            \,\middle|\,
            \mathcal U_A
        \right]
        \le
        C e^{-(1+\kappa) T}
        e^{-\lambda_0W(A)}e^{c\beta h},
\end{aligned}
\]
by \eqref{eq:MA-bound}.
For the terminal branch,
\[
\begin{aligned}
        \mathbb P(E_{{\rm fail},A}^*(h)\mid\mathcal U_A)
        &=
        e^{-(1+\kappa) T}
        \mathbb E
        \left[
            e^{(1+\kappa) T}
            \mathbf 1_{E_{{\rm fail},A}^*(h)}
            \,\middle|\,
            \mathcal U_A
        \right]  \le
        C e^{-(1+\kappa) T}
        e^{-\lambda_0W(A)}e^{c\beta h}.
\end{aligned}
\]
Since \(E_A^*(h)=E_{{\rm res},A}^*(h)\cup E_{{\rm fail},A}^*(h)\),
\[
        \mathbb P(E_A^*(h)\mid\mathcal U_A)
        \le
        C e^{-(1+\kappa)T}e^{-\lambda_0W(A)}e^{c\beta h}.
\]
Choose
\[
        \lambda_0:=\lambda_{\rm SR}+\log q+\log c_{\rm bl}^{-1}+1.
\]
Using \(W(A)\ge |A|\), the preceding display and the bound before this
paragraph give
\[
\begin{aligned}
    q^{|A|}
    \mathbb P
    \left(
        A\in{\rm SR}
        \,\middle|\,
        \mathcal H_A^-=X,\mathsf T_A
    \right)
    &\le
    C
    e^{-(1+\kappa)T}
    e^{-\lambda_{\rm SR} W(A)}
    e^{c\beta h}.
\end{aligned}
\]
Renaming the constants as \(C_{\rm sr}\) and \(c_{\rm sr}\) proves the claim.





\end{proof}

\end{document}
